\documentclass[11pt]{article}
\usepackage[a4paper,margin=25mm]{geometry}
\usepackage{amsmath,amssymb,amsthm,mathtools,microtype}
\usepackage[T1]{fontenc}\usepackage{lmodern}
\usepackage[colorlinks=true,linkcolor=blue,urlcolor=blue]{hyperref}
\newtheorem{theorem}{Theorem}[section]
\newtheorem{lemma}[theorem]{Lemma}
\newtheorem{corollary}[theorem]{Corollary}
\newcommand{\F}{\mathbb F}
\newcommand{\SV}{\operatorname{SV}}
\newcommand{\wt}{\operatorname{wt}}
\newcommand{\pdeg}{\operatorname{pdeg}}
\newcommand{\Span}{\operatorname{span}}
\newcommand{\rank}{\operatorname{rank}}
\newcommand{\GL}{\operatorname{GL}}
\title{Affine-Orbit Hitting for Invertible-Local ROABPs\\
in Every Characteristic}
\author{Yilin Liu \and GPT-6 Astra Pro}\date{26 October 2026}
\begin{document}\maketitle

\begin{abstract}
We prove an exponential-in-width bound on the characteristic digit degree
of the entire first homogeneous scalar component of an invertible-local
read-once oblivious algebraic branching program after a separating curve
shift. For each fixed support of size $k$, an exterior pairing bounds the
number of leading monomials by $\binom{w-2}{k-1}$. Full common-translation
invariance forces whole-digit transfers among monomials of maximal digit
degree. A finite-set spanning-tree argument then gives
$\pdeg\le b(w):=\max_k\{\binom{w-2}{k-1}+k-1\}$ for $w\ge2$.
Consequently a curve and $b(w)$ Shpilka--Volkovich pairs hit unknown
invertible affine orbits in every characteristic. Univariate rank-one
splitting permits local generic ranks one or $w$. A self-contained digit
polarization argument gives degree-independent direction grids and
$n^{2b(w)}(\rho D+1)(D+1)^{b(w)}$ evaluations for scalar total degree
at most $D$, where $\rho\le2n-1$.
\end{abstract}

\noindent\textbf{Review status.}
This manuscript and its preceding structural notes have independent
model-assisted internal reconstructions and exact-source audits.
The mathematical text is frozen. No external peer review
or publication-priority claim is asserted.

\section{Statement and precise model}

Let $n,w,D\ge1$. We consider a nonzero polynomial
\begin{equation}\label{eq:model}
 f(X)=uA_1(\ell_1(X)+c_1)\cdots A_n(\ell_n(X)+c_n)v,
 \qquad \deg f\le D,
\end{equation}
over a field $\F$. Each $A_i$ is a $w$-by-$w$ univariate polynomial
matrix of generic rank either one or $w$; $u,v$ are constant endpoints;
and $\ell_1,\ldots,\ell_n$ are the coordinate forms of an invertible
linear transformation. The representation, its order, and the affine
transformation are unknown. There is no characteristic or local-degree
restriction. A representation with local entry degrees at most $\delta$
has $D\le n\delta$. A known latent scalar individual-degree bound
$\delta$ also implies $D\le n\delta$, without changing the promised
representation.

Put
\begin{equation}\label{eq:b}
 b(1)=0,\qquad
 b(w)=\max_{1\le k\le w-1}
       \left\{\binom{w-2}{k-1}+k-1\right\}\quad(w\ge2).
\end{equation}
For widths $2,\ldots,8$ the values are $1,2,3,5,8,13,23$.
For $w\ge2$, the estimate $b(w)\le2^{w-2}+w-2$ suffices for the asymptotic statement.
No optimality of this width-dependent constant is asserted.

Choose $n$ distinct interpolation nodes $a_i$ in a working field and
their Lagrange polynomials $\lambda_i$, of degree at most $n-1$.
For $k\ge0$ define
\[
 \SV_{k,i}=\sum_{j=1}^k z_j\lambda_i(y_j).
\]
In characteristic zero let $e_i=i$. In characteristic $p>0$, let $e_i$
be the $i$-th positive integer not divisible by $p$. Set
\[
 \gamma_i(t)=t^{e_i},\qquad \rho=e_n,\qquad n\le\rho\le2n-1.
\]
The coordinate derivatives of $\gamma$ are linearly independent.

\begin{theorem}\label{thm:main}
Every nonzero polynomial \eqref{eq:model}, over every characteristic,
is hit by
\begin{equation}\label{eq:map}
 \mathcal H_i(t,z,y)=t^{e_i}
          +\sum_{j=1}^{b(w)}z_j\lambda_i(y_j).
\end{equation}
The standard map $\SV_{b(w)+1}$ also hits this class.
For \eqref{eq:map}, any Cartesian grids with
\[
 |T|=\rho D+1,\qquad |Z_j|=D+1,\qquad |Y_j|=n^2
\]
give a hitting set of size at most
\begin{equation}\label{eq:count}
 n^{2b(w)}(\rho D+1)(D+1)^{b(w)}.
\end{equation}
A working field of size at least $\max\{n^2,\rho D+1\}$ supplies
all these grids and the interpolation nodes. Extension-field oracle
evaluations are allowed when the base field is smaller.
If $D=n\delta$ with $n,\delta\ge1$, the count is
$O_w(n^{3b(w)+2}\delta^{b(w)+1})$.
Constant polynomials, including $n=0$ or $D=0$, need one evaluation.
\end{theorem}

The theorem permits rank-one bottlenecks, but does not include local
generic ranks strictly between one and $w$. It concerns invertible
affine transformations, not singular projections or orbit closures.
The proof bounds the \emph{whole} first homogeneous component. It does
not infer affine hitting merely from the existence of one coefficient
of low digit degree.

\section{Identity normalization and full translation}

In characteristic $p>0$, write $\wt_p(a)$ for the sum of the base-$p$
digits of a nonnegative integer $a$. The digit degree of $Y^\alpha$ is
$\sum_i\wt_p(\alpha_i)$, and that of a nonzero polynomial is its
maximum monomial digit degree. In characteristic zero define
$\wt_0(a)=a$, so digit degree means ordinary total degree.

We first allow generically invertible rational univariate factors
$H_i(T_i)$ of width $w$, over an arbitrary coefficient field $F_0$.
Let $u,v$ be constant nonzero endpoints and
\[
 f(T)=uH_1(T_1)\cdots H_n(T_n)v\ne0.
\]
Let $C_i(t)\in F_0(t)$ satisfy $C_i'(t)\ne0$ for every $i$.
Each $C_i$ is transcendental over $F_0$, so each nonzero local
determinant and univariate denominator remains nonzero on this curve.
Thus $f(C+Y)$ is regular at $Y=0$ over $E=F_0(t)$.
No perfection assumption on either coefficient field is used.

If $f(C)\ne0$, its first component is a nonzero constant. Otherwise put
\[
 P_j=H_1(C_1)\cdots H_j(C_j),\qquad
 B_i(Y_i)=P_{i-1}H_i(C_i+Y_i)P_i^{-1},\qquad v_*=P_nv.
\]
Then $B_i(0)=I_w$, $u,v_*$ are nonzero, $uv_*=0$, and
\begin{equation}\label{eq:identity}
 g(Y)=f(C+Y)=uB_1(Y_1)\cdots B_n(Y_n)v_*.
\end{equation}
Let $m>0$ be the first nonzero ordinary total degree of this formal
series, and let $H=g_{[m]}$ be its homogeneous component of degree $m$.

\begin{lemma}\label{lem:translation}
The leading component satisfies the polynomial identity
\[
 H_t(Y+hC'(t))=H_t(Y).
\]
After the diagonal change $Y_i=C_i'(t)Z_i$, it satisfies
$H(Z+h\mathbf1)=H(Z)$. This change preserves support, digit degree,
and the number of nonzero monomials on every support.
\end{lemma}
\begin{proof}
Expand $f(C(t+h)+Y)$ as a formal series in the combined variables
$(Y,h)$. Expanding first at the curve parameter $t+h$, every component
of residual $Y$-degree below $m$ vanishes identically as a rational
function of the parameter. Its total-degree-$m$ part is therefore
$H_t(Y)$; changing a coefficient from $t$ to $t+h$ affects only higher
combined degree. Expanding instead at $C(t)$ gives
$f(C(t)+Y+\Delta)$ with
$\Delta_i=hC_i'(t)+O(h^2)$. Its total-degree-$m$ part is
$H_t(Y+hC')$. Comparing the expansions proves the full identity.
All $C_i'$ are nonzero, which proves the remaining assertions.
\end{proof}

The full translation identity, rather than only its first derivative,
is essential in positive characteristic.

\section{An exterior bound for each fixed support}

The following lemma applies to any identity-constant word
\eqref{eq:identity}, even without its origin on a curve. Write
$B_{i,a}=[Y_i^a]B_i$.

\begin{lemma}\label{lem:exterior}
Let $w\ge2$. For the first nonzero homogeneous component $H$ of an identity-constant
width-$w$ word with $uv_*=0$, the number $N_S$ of nonzero monomials on
a fixed support $S$ of size $k$ satisfies
\[
 N_S\le\binom{w-2}{k-1}.
\]
In particular every nonzero leading monomial has support at most $w-1$.
\end{lemma}
\begin{proof}
The case $k=1$ follows from homogeneity. Assume $k\ge2$ and write
$S=\{i_1<\cdots<i_k\}$. For a nonzero leading monomial $Y^\alpha$
on $S$, write its positive exponents as $\alpha_1,\ldots,\alpha_k$,
whose sum is $m$, and define
\[
 r_s^\alpha=uB_{i_1,\alpha_1}\cdots B_{i_s,\alpha_s},\qquad
 c_t^\alpha=B_{i_{t+1},\alpha_{t+1}}\cdots B_{i_k,\alpha_k}v_* .
\]
For $1\le s,t\le k-1$, the rows $r_s^\alpha$ lie in
$V=\ker(v_*)$ and the columns $c_t^\alpha$ lie in $W=\ker(u)$:
their scalar contractions are coefficients of ordinary degree below
$m$. The pairing induced on $V/Eu$ and $W/Ev_*$ is nondegenerate and
has dimension $w-2$. Indeed the annihilator of $W$ in the row space is
$Eu$, and dually the annihilator of $V$ is $Ev_*$.

Take the exterior products of the $k-1$ prefix rows and the $k-1$
suffix columns in these quotient spaces. Their exterior pairing is
the determinant of $(r_s^\alpha c_t^\beta)_{s,t=1}^{k-1}$.
For $\alpha=\beta$, entries with $s<t$ vanish: their variable sets
are disjoint and their total degree is below $m$. Every diagonal entry
is the nonzero full coefficient $h_\alpha=[Y^\alpha]H$. Thus
\begin{equation}\label{eq:diagonal}
 \det(r_s^\alpha c_t^\alpha)=h_\alpha^{k-1}\ne0.
\end{equation}

Order the exponent vectors on $S$ lexicographically. Suppose
$\alpha<\beta$, with first difference at position $h$. Then
$\alpha_h<\beta_h$ and $h\le k-1$, since their total degrees agree.
Let $a_s=\sum_{j\le s}\alpha_j$ and
$b_t=\sum_{j\le t}\beta_j$. For $s\le h$ and $t\ge h$,
\[
 a_s\le a_h<b_h\le b_t.
\]
Here $s\le t$, so the prefix and suffix use disjoint variable sets;
their product is a genuine scalar coefficient of total degree
$a_s+m-b_t<m$. It vanishes. The first $h$ rows of the pairing matrix
are consequently supported on only $h-1$ columns, and its determinant
is zero.

The matrix of exterior pairings over all leading monomials on $S$ is
triangular with nonzero diagonal \eqref{eq:diagonal}. Its rank is
$N_S$ and is at most the dimension of the exterior space,
$\binom{w-2}{k-1}$. If $k\ge w$, this dimension is zero, so no
such monomial exists. Entries with $s>t$ may repeat variable positions;
no coefficient interpretation or vanishing assertion for those entries
was used.
\end{proof}

\section{Digit transfers and a spanning-tree count}

\begin{lemma}\label{lem:sets}
Let a finite connected undirected graph have vertex labels
$d(v)=(d_1(v),\ldots,d_k(v))$ in positive integers, all with sum $K$.
Every edge transfers one unit from one coordinate to another. Suppose
that for every vertex and every coordinate $i$ with $d_i>1$, there is
an adjacent vertex whose $i$-th coordinate is $d_i-1$.
Then the graph has at least $K-k+1$ vertices. Labels need not be distinct.
\end{lemma}
\begin{proof}
Put $R=K-k$ and associate to vertex $v$ the $R$-element set
\[
 S_v=\{(i,r):1\le r\le d_i(v)-1\}.
\]
Adjacent sets exchange one element. For every fixed coordinate, repeated
decreasing steps reach weight one within the connected component.
Hence $\bigcap_v S_v=\varnothing$.

Order the vertices along a rooted spanning tree, with each parent
preceding its children. The initial intersection has $R$ elements.
Before adding a child, the current intersection is contained in its
parent's set. The child removes at most one element of that set, so
each added vertex decreases the intersection by at most one. An empty
final intersection requires at least $R$ additions, proving the claim.
The argument also covers $R=0$ and repeated labels.
\end{proof}

\begin{lemma}\label{lem:transfer}
Let $H$ be a nonzero homogeneous polynomial with
$H(Z+h\mathbf1)=H(Z)$, and let $K$ be its maximal digit degree.
Choose a support $S$ of maximum cardinality $k$ among its nonzero
digit-degree-$K$ monomials. If $N$ is the number of such monomials on
this same support, then
\[
 K\le N+k-1.
\]
\end{lemma}
\begin{proof}
Let $Z^\alpha$ be any of these monomials and put
$d_i=\wt_p(\alpha_i)$, using ordinary exponents in characteristic zero.
If $d_i>1$, choose a digit power $q=p^e$ that occurs in $\alpha_i$;
in characteristic zero take $q=1$. Then $\alpha_i-q>0$ and its digit
weight is $d_i-1$.

In the translation identity the monomial $Z^\alpha$ contributes
nontrivially to
\[
 Z^\gamma h^q,\qquad \gamma=\alpha-qe_i.
\]
This coefficient must cancel. Any contributing source has exponent
$\beta=\gamma+\delta$, with $\delta_j\ge0$ and
$\sum_j\delta_j=q$. In positive characteristic, the nonzero binomial
coefficient requires that every digit of $\gamma_j$ be a subdigit of
$\beta_j$. This follows directly by expanding
$(Z+h)^a=\prod_e(Z^{p^e}+h^{p^e})^{a_e}$.
There is therefore no borrow in the subtraction, and
\begin{equation}\label{eq:weight}
 \wt_p(\beta)=\wt_p(\gamma)+\sum_j\wt_p(\delta_j)
              =K-1+\sum_j\wt_p(\delta_j).
\end{equation}
The same equation holds with ordinary weights in characteristic zero.
Maximality of $K$ forces the last sum to be one. Since
$\sum_j\delta_j=q$, exactly one increment is nonzero and it equals
$q$. A distinct nonzero cancelling source thus has the form
\begin{equation}\label{eq:transfer}
 \beta=\alpha-qe_i+qe_j,\qquad j\ne i,
\end{equation}
with no carry in the addition at $j$ and the same digit degree $K$.
If $j$ were outside $S$, its support would have size $k+1$, because
the decreased coordinate remains positive. This contradicts the choice
of $k$. The partner lies on the same support.

Form the graph of all $N$ monomials, with edges for all such permitted
whole-digit transfers. An edge is reversible: its destination has digit
weight at least two, removing the transferred digit has no borrow, and
adding it back restores the original exponent without a carry.
Each connected component has the decreasing-step property of
Lemma~\ref{lem:sets} for every coordinate. That lemma gives at least
$K-k+1$ vertices in any component, and hence $N\ge K-k+1$.
Only the existence of a cancelling partner was used; there is no claim
that every destination or every allocation of digit copies must occur.
\end{proof}

\begin{theorem}\label{thm:component}
For the rational invertible-local word and curve of Section 2, the
entire first nonzero homogeneous component of $f(C(t)+Y)$ has digit
degree at most $b(w)$. In width one, $f(C(t))\ne0$.
\end{theorem}
\begin{proof}
If the constant component is nonzero there is nothing to prove.
Otherwise use Lemma~\ref{lem:translation} to make common translation
invariance hold by a nonzero diagonal scaling. This preserves the
fixed-support counts in Lemma~\ref{lem:exterior}. For a maximal-digit-
degree monomial on a largest such support $S$, Lemma~\ref{lem:transfer}
and the exterior bound give
\[
 K\le N+k-1\le\binom{w-2}{k-1}+k-1\le b(w).
\]
Here $1\le k\le w-1$. In width one the scalar is a product of
nonzero univariate rational functions and nonzero constant endpoints;
each factor stays nonzero at its transcendental $C_i$.
\end{proof}

\section{Affine changes and rank-one splitting}

Digit degree cannot increase under linear substitution. In positive
characteristic this follows by expanding each exponent into base-$p$
digits and using Frobenius on every digit copy; subsequent collisions
can only reduce the resulting digit weight. An invertible linear
substitution also preserves nonzeroness. A polynomial of digit degree
at most $k$ has every monomial supported on at most $k$ variables.

The image of $\SV_k$ contains every coordinate subspace of dimension
at most $k$: specialize its direction variables to the corresponding
interpolation nodes, and unused amplitudes to zero. Consequently it
hits every nonzero polynomial having a nonzero monomial supported on
at most $k$ variables. Indeed restriction to that monomial's coordinate
set retains its coefficient.

For \eqref{eq:model} with all local factors invertible, set
$C_i(t)=c_i+\ell_i(\gamma(t))$. Independent coordinate derivatives
ensure $C_i'\ne0$. Theorem~\ref{thm:component} applies. The residual
invertible linear change takes the entire leading component to a
nonzero polynomial of digit degree at most $b(w)$, which is hit by
$\SV_{b(w)}$. Higher homogeneous components have larger total degree
in the fresh amplitudes and cannot cancel it. This proves
$f(\gamma+\SV_{b(w)})\ne0$ for invertible local factors.
The argument applies to any subset of independent latent forms by
extending them to a full basis of the original coordinate space.

For completeness we give the rank-one interface explicitly. Let
$R_j(T_j)$, $1\le j\le s$, be the rank-one factors in their order,
and let $G_0,\ldots,G_s$ denote the intervening invertible words.
Choose a nonzero univariate pivot $h_j=(R_j)_{a_jb_j}$ and put
\[
 a_j^*=R_je_{b_j},\qquad b_j^*=e_{a_j}^{T}R_j,
 \qquad R_j=\frac{a_j^*b_j^*}{h_j}.
\]
The scalar has the exact rational factorization
\begin{equation}\label{eq:split}
 f=\frac{(uG_0a_1^*)
       (b_1^*G_1a_2^*)\cdots
       (b_{s-1}^*G_{s-1}a_s^*)(b_s^*G_sv)}
          {h_1\cdots h_s}.
\end{equation}
Every numerator is nonzero because $f\ne0$, and is internally read
once. Adjacent numerators may share a rank-one boundary variable.

A nonzero univariate row or column can be completed to an invertible
univariate matrix by adjoining constant standard-basis rows or columns:
choose one nonzero coordinate for its determinant. Make these choices
over the original coefficient field, before evaluating the curve.
Thus every numerator in \eqref{eq:split} is an invertible-local scalar
word with constant endpoints, to which the preceding argument applies.
The completion determinants and the pivots $h_j$ remain nonzero at
the curve because their latent arguments are transcendental. Their
shifted denominators are units at residual zero. The same map preserves
all numerator factors, and their product remains nonzero even though
different factors can share variables. There are at most $s+1\le n+1$
numerators. If $s=0$, use the one original scalar word instead.

The structural argument allowed rational univariate factors precisely
so that regular univariate endpoint completions or pivot divisions do
not require a change of coefficient field depending on residual
variables. All such rational functions are defined over the original
field; the only common coefficient extension in the proof is $F_0(t)$.

\section{Digit polarization and finite direction grids}

We include the direction-grid proof to make the query count independent
of any auxiliary addendum.

\begin{lemma}\label{lem:polar}
Let $0\ne H$ be homogeneous of positive ordinary degree over any
field $E$, with maximal digit degree $K$. There are $K$ pure digit
powers $q_1,\ldots,q_K$ (all one in characteristic zero) such that
\[
 \mathcal P(V_1,\ldots,V_K)=
 [z_1^{q_1}\cdots z_K^{q_K}]
 H\left(\sum_{j=1}^K z_jV_j\right)
\]
is a nonzero tensor of the form
$\sum c_{i_1,\ldots,i_K}\prod_j(V_j)_{i_j}^{q_j}$.
It remains nonzero after any invertible linear change in every vector
argument induced by the same invertible coordinate map.
\end{lemma}
\begin{proof}
Choose a nonzero monomial $h_\alpha Y^\alpha$ of digit degree $K$.
In characteristic $p$, write $\alpha_i=\sum_e d_{i,e}p^e$,
$0\le d_{i,e}<p$, and list its $K$ digit copies in any order.
Assign one amplitude to each copy, recording its coordinate $i_j$
and power $q_j$. A source monomial contributes to the target amplitude
monomial only if all $K$ amplitudes are occupied. Every source has at
most $K$ digit copies, so exactly one copy must occupy each amplitude.
This proves the tensor form. The coefficient of the particular vector
monomial $\prod_j(V_j)_{i_j}^{q_j}$ is
\[
 h_\alpha\prod_{i,e}d_{i,e}!\ne0.
\]
Its coordinate/exponent list determines $\alpha$, so a different
source cannot cancel this coefficient. In characteristic zero the same
argument uses $K=\deg H$, powers one, and coefficient
$h_\alpha\prod_i\alpha_i!\ne0$.
Coefficient extraction commutes with an invertible coordinate change,
and replacing every $V_j$ by $LV_j$ is invertible. Equivalently, each
Frobenius-twisted matrix has determinant $(\det L)^{q_j}\ne0$.
No total factorial or perfection assumption is required.
\end{proof}

The polynomial in Lemma~\ref{lem:polar} remains nonzero when
$V_j=(\lambda_1(y_j),\ldots,\lambda_n(y_j))$. Each family
$\{\lambda_i(y_j)^{q_j}\}_i$ is independent over every coefficient
extension, as evaluation at the interpolation nodes gives the identity
matrix.

\begin{lemma}\label{lem:directions}
Suppose $M$ nonzero polynomials in $k$ direction variables have, in
each variable $y_j$, coefficients in the span of
$\{\lambda_i(y_j)^{q_j}\}_i$, where $q_j$ is a pure digit power;
the power lists may differ between polynomials. A polynomial constant
in a direction is allowed. Any $k$ grids of size $M(n-1)+1$ contain
a tuple preserving all $M$ polynomials as nonzero over the coefficient
field.
\end{lemma}
\begin{proof}
Choose directions successively. In each currently nonzero polynomial,
select a nonzero coefficient after expanding in the remaining
direction variables. As a polynomial in the current direction it is
$\sum_i b_i\lambda_i(y)^{q}$. Over an algebraic closure of the
coefficient field it is a $q$-th power of a nonzero polynomial of
degree at most $n-1$. Hence it has at most $n-1$ distinct roots.
In characteristic zero $q=1$ and the same bound is immediate.
Avoid the union of at most $M(n-1)$ bad choices and continue.
Constant directions have no bad choices. Root extraction is used only
in the proof; the algorithm searches the specified Cartesian grids.
\end{proof}

Apply these lemmas over $E=\F(t)$ to the at most $M\le n+1$
numerators in \eqref{eq:split}. A numerator already nonzero at zero
uses its nonzero constant amplitude coefficient. For each other one,
apply polarization to its entire first homogeneous component, of digit
degree $K\le b(w)$, using the first $K$ amplitude variables and
exponent zero in the remaining ones. The target amplitude degree is
the ordinary degree of that component, so no higher component changes
this coefficient. Direction polynomials for different numerators may
use different powers. Lemma~\ref{lem:directions} simultaneously
preserves all of them on grids of size
\[
 M(n-1)+1\le(n+1)(n-1)+1=n^2.
\]
Every auxiliary denominator stays a unit at amplitude zero regardless
of those directions. Thus the original $f$ remains a nonzero polynomial
in $t,z_1,\ldots,z_{b(w)}$ after one such direction choice.

Its individual degrees in these remaining variables are at most
$\rho D,D,\ldots,D$. Standard individual-degree interpolation on
$\rho D+1,D+1,\ldots,D+1$ points gives a nonzero evaluation.
This proves \eqref{eq:count} and its field-size statement. The search
evaluates only the original polynomial composed with the polynomial
map \eqref{eq:map}; numerical grid points need not avoid any auxiliary
rational pivot. In width one, no directions or fresh amplitudes are
needed and the same argument uses only the curve grid.

For comparison, ordinary degree interpolation in all seeds gives the
basic rectangular count
\[
 (\rho D+1)(D+1)^{b(w)}\bigl((n-1)D+1\bigr)^{b(w)}.
\]
The improved count replaces its direction grids by $n^2$ elements
without imposing any bound on $D$ relative to the characteristic.

\section{The pure Shpilka--Volkovich map}

Use the first pair of $\SV_{b(w)+1}$ as a proof-only separating
curve. In characteristic $p>0$, specialize
$z_0=t$, $y_0=t^p$. The coordinate derivatives of
$t\lambda_i(t^p)$ are $\lambda_i(t^p)$ and are independent.
In characteristic zero, specialize $z_0=t$, $y_0=t$. The derivatives
are $\lambda_i(t)+t\lambda_i'(t)$; the map
$q(t)\mapsto q(t)+tq'(t)$ is invertible on polynomials of degree
below $n$, since it multiplies $t^a$ by $a+1$.
The remaining $b(w)$ pairs are fresh, so the already proved argument
shows that this specialization has nonzero image. The unspecialized
$\SV_{b(w)+1}$ image is consequently nonzero as well.

A basic finite grid for this pure map has $(D+1)$ values per amplitude
and $((n-1)D+1)$ per direction, for at most
\[
 \bigl((D+1)((n-1)D+1)\bigr)^{b(w)+1}
 \le(nD+1)^{2b(w)+2}
\]
evaluations. This bound uses the actual polynomial map, not the
possibly higher-degree proof-only curve. In width one it gives
$\SV_1$, consistent with the curve-only argument.

\section{Limits of the conclusion}

The exponent depends only on width and is at most exponential in that
width. The proof does not establish the smaller bound $w-1$ in every
characteristic. Its combinatorial lower bound is $N\ge K-k+1$,
not $N\ge\binom{K-1}{k-1}$; the latter assertion is false in general
even under full translation invariance. Intermediate singular local
ranks require additional arguments and are outside the present model.
All finite-grid degree bounds refer to the original scalar polynomial,
so a known scalar degree promise can be used without truncating or
otherwise changing the rank pattern of the representation.

\end{document}
